\documentclass[a4paper,12pt]{article}

\usepackage[T2A]{fontenc}
\usepackage[utf8]{inputenc}
\usepackage[russian]{babel}

\usepackage{amsmath, amssymb, amsthm, mathtools}
\usepackage{helvet}
\renewcommand{\familydefault}{\sfdefault}
\usepackage{icomma}
\usepackage[left=18mm,right=18mm,top=18mm,bottom=20mm]{geometry}
\usepackage{microtype}
\usepackage{tikz}
\usepackage{tkz-euclide}
\usepackage{pgfplots}
\pgfplotsset{compat=1.18}
\usepackage{tabularx, booktabs, longtable}
\usepackage{enumitem}
\usepackage{hyperref}
\usepackage{parskip}
\usepackage{fancyhdr}
\usepackage{setspace}
\usepackage{xcolor}

\definecolor{accent}{HTML}{008080}
\definecolor{darkaccent}{HTML}{004D4D}
\definecolor{textgray}{HTML}{333333}
\definecolor{softgray}{HTML}{F4F7F7}

\providecommand{\tg}{\mathop{\mathrm{tg}}\nolimits}
\providecommand{\ctg}{\mathop{\mathrm{ctg}}\nolimits}
\providecommand{\arctg}{\mathop{\mathrm{arctg}}\nolimits}

\hypersetup{
    colorlinks=true,
    linkcolor=darkaccent,
    urlcolor=darkaccent
}

\pagestyle{fancy}
\fancyhf{}
\lhead{\textcolor{textgray}{Логарифмические неравенства}}
\rhead{\textcolor{textgray}{ЕГЭ}}
\cfoot{\thepage}
\renewcommand{\headrulewidth}{0.4pt}
\renewcommand{\headrule}{\hbox to\headwidth{\color{accent}\leaders\hrule height \headrulewidth\hfill}}

\onehalfspacing
\setlength{\parindent}{0pt}
\setlength{\emergencystretch}{2em}
\setlist[itemize]{leftmargin=1.2em, itemsep=0.15em, topsep=0.2em}
\setlist[enumerate]{leftmargin=1.4em, itemsep=0.35em, topsep=0.3em}
\allowdisplaybreaks

\newcommand{\sectionline}{\vspace{0.35em}{\color{accent}\hrule height 0.6pt}\vspace{0.65em}}
\newcommand{\important}[1]{\textcolor{darkaccent}{\textbf{#1}}}

\begin{document}

\begin{titlepage}
\thispagestyle{empty}
\begin{center}
\vspace*{2.5cm}

{\Large \textcolor{accent}{\textbf{Чек-лист}}}

\vspace{0.7cm}

{\huge \textbf{Логарифмические неравенства}}\\[0.25cm]
{\Large \textbf{ОДЗ, автоматические ограничения, модули и гибридные задачи}}

\vspace{1.3cm}

{\large Лёвин Артём Александрович эксклюзивно для Николь}

\vspace{0.6cm}

{\large \today}

\vfill

\begin{tikzpicture}[scale=1]
\draw[accent, line width=1.1pt]
    (-5.8,0)
    .. controls (-3.7,0.8) and (-2.1,-0.8) .. (0,0)
    .. controls (2.1,0.8) and (3.7,-0.8) .. (5.8,0);
\draw[darkaccent, line width=0.45pt]
    (-4.4,-0.45)
    .. controls (-2.7,0.2) and (-1.3,-0.55) .. (0,-0.25)
    .. controls (1.3,0.05) and (2.7,-0.65) .. (4.4,-0.05);
\end{tikzpicture}

\vspace*{1.2cm}
\end{center}
\end{titlepage}

\section*{\textcolor{darkaccent}{1. Главная идея занятия}}
\sectionline

Логарифмические неравенства требуют не только знания формул, но и внимательного учёта ограничений. В задачах ЕГЭ часто встречаются неравенства, где одновременно используются:

\begin{itemize}
    \item логарифмы;
    \item дроби;
    \item модули;
    \item показательные выражения;
    \item автоматические ограничения, которые можно не решать отдельно.
\end{itemize}

\important{Главная цель:} научиться видеть, какие ограничения действительно нужно решать, а какие уже выполняются автоматически.

\section*{\textcolor{darkaccent}{2. Базовые условия для логарифма}}
\sectionline

Логарифм
\[
\log_a f(x)
\]
имеет смысл только при выполнении условий:
\[
a>0,\qquad a\ne1,\qquad f(x)>0.
\]

Если основание задано числом, например $2$, $4$, $\dfrac13$, то условия на основание обычно не пишут, потому что они очевидны:
\[
2>0,\qquad 2\ne1,
\]
\[
\frac13>0,\qquad \frac13\ne1.
\]

\important{В решении нужно контролировать именно подлогарифмическое выражение:}
\[
f(x)>0.
\]

\section*{\textcolor{darkaccent}{3. Как отбрасывать логарифмы}}
\sectionline

Пусть
\[
\log_a f(x) \quad \text{и} \quad \log_a g(x)
\]
имеют одно и то же основание.

\subsection*{\textcolor{accent}{Если основание больше единицы}}

При $a>1$ логарифмическая функция возрастает, поэтому знак неравенства сохраняется:
\[
\log_a f(x)>\log_a g(x)
\quad\Longleftrightarrow\quad
f(x)>g(x).
\]

Но ограничения по-прежнему важны:
\[
f(x)>0,\qquad g(x)>0.
\]

Иногда одно из них выполняется автоматически. Например, если
\[
g(x)>0,\qquad f(x)>g(x),
\]
то
\[
f(x)>0
\]
уже следует из этих двух условий.

\subsection*{\textcolor{accent}{Если основание между нулём и единицей}}

При $0<a<1$ логарифмическая функция убывает, поэтому знак неравенства меняется:
\[
\log_a f(x)>\log_a g(x)
\quad\Longleftrightarrow\quad
f(x)<g(x).
\]

Ограничения остаются такими же:
\[
f(x)>0,\qquad g(x)>0.
\]

Знак меняется только при сравнении аргументов, но не в условиях существования логарифмов.

\section*{\textcolor{darkaccent}{4. Какие ограничения можно не писать}}
\sectionline

\subsection*{\textcolor{accent}{Правило ``меньшей стороны''}}

После перехода от логарифмов к аргументам достаточно оставить положительность той стороны, которая оказалась \important{не больше} другой.

\begin{table}[h]
\caption{Автоматические ограничения}
\centering
\renewcommand{\arraystretch}{1.45}
\begin{tabularx}{\linewidth}{>{\bfseries}p{0.23\linewidth}p{0.31\linewidth}X}
\toprule
Основание & Переход & Что достаточно оставить \\
\midrule
$a>1$ &
$\log_a f(x)>\log_a g(x)\Rightarrow f(x)>g(x)$ &
достаточно $g(x)>0$, потому что $f(x)>g(x)>0$ \\

$a>1$ &
$\log_a f(x)<\log_a g(x)\Rightarrow f(x)<g(x)$ &
достаточно $f(x)>0$, потому что $g(x)>f(x)>0$ \\

$0<a<1$ &
$\log_a f(x)>\log_a g(x)\Rightarrow f(x)<g(x)$ &
достаточно $f(x)>0$, потому что $g(x)>f(x)>0$ \\

$0<a<1$ &
$\log_a f(x)<\log_a g(x)\Rightarrow f(x)>g(x)$ &
достаточно $g(x)>0$, потому что $f(x)>g(x)>0$ \\
\bottomrule
\end{tabularx}
\end{table}

То же правило работает для нестрогих знаков $\le$ и $\ge$, если сохранённое ограничение строгое:
\[
g(x)>0,\qquad f(x)\ge g(x)
\quad\Rightarrow\quad
f(x)>0.
\]

\section*{\textcolor{darkaccent}{5. Константу удобно записывать как логарифм}}
\sectionline

Чтобы сравнить логарифм с числом, число часто переводят в логарифмический вид:
\[
c=\log_a a^c.
\]

\subsection*{\textcolor{accent}{Примеры}}

\[
1=\log_2 2,
\qquad
2=\log_2 4,
\qquad
\frac12=\log_4 2,
\qquad
-3=\log_{\frac12}8.
\]

Например:
\[
\log_4 F(x)\ge\frac12
\]
можно переписать как
\[
\log_4 F(x)\ge \log_4 2.
\]

Так как $4>1$, получаем:
\[
F(x)\ge2.
\]

Ограничение $F(x)>0$ здесь выполняется автоматически, потому что
\[
F(x)\ge2>0.
\]

\section*{\textcolor{darkaccent}{6. Общий алгоритм решения логарифмического неравенства}}
\sectionline

\begin{enumerate}
    \item Найдите основание логарифма.
    \item Проверьте, возрастает или убывает логарифмическая функция:
    \[
    a>1 \quad \text{или} \quad 0<a<1.
    \]
    \item Запишите ОДЗ для исходного выражения.
    \item Если справа число, представьте его в виде логарифма:
    \[
    c=\log_a a^c.
    \]
    \item Отбросьте логарифмы:
    \begin{itemize}
        \item при $a>1$ знак сохраняется;
        \item при $0<a<1$ знак меняется.
    \end{itemize}
    \item Проверьте, можно ли часть ограничений не решать отдельно.
    \item Решите полученную систему или совокупность.
    \item Проверьте, не нарушены ли знаменатели, модули и исходная ОДЗ.
\end{enumerate}

\section*{\textcolor{darkaccent}{7. Мини-примеры равносильных переходов}}
\sectionline

\subsection*{\textcolor{accent}{Пример 1. Ограничение выполняется автоматически}}

\[
\log_2\frac{x^2-3x}{x+1}\ge0.
\]

Так как
\[
0=\log_2 1,
\]
получаем:
\[
\log_2\frac{x^2-3x}{x+1}\ge\log_2 1.
\]

Основание $2>1$, значит:
\[
\frac{x^2-3x}{x+1}\ge1.
\]

Ограничение
\[
\frac{x^2-3x}{x+1}>0
\]
можно не решать отдельно, потому что из
\[
\frac{x^2-3x}{x+1}\ge1
\]
сразу следует положительность дроби.

\subsection*{\textcolor{accent}{Пример 2. Ограничение нельзя отбросить}}

\[
\log_2\frac{x^2-3x}{x+1}<0.
\]

Получаем:
\[
\log_2\frac{x^2-3x}{x+1}<\log_2 1.
\]

Так как $2>1$:
\[
\frac{x^2-3x}{x+1}<1.
\]

Но это не гарантирует положительность дроби. Поэтому нужно решать систему:
\[
\begin{cases}
\dfrac{x^2-3x}{x+1}>0,\\[0.6em]
\dfrac{x^2-3x}{x+1}<1.
\end{cases}
\]

\section*{\textcolor{darkaccent}{8. Логарифмы с модулем}}
\sectionline

Модуль часто позволяет упростить ОДЗ, но не всегда.

\subsection*{\textcolor{accent}{Если модуль стоит на всём аргументе}}

Пусть
\[
\log_a |F(x)|
\]
входит в неравенство. Тогда условие существования:
\[
|F(x)|>0.
\]

Так как модуль всегда неотрицателен, это равносильно условию:
\[
F(x)\ne0.
\]

\subsection*{\textcolor{accent}{Пример}}

\[
\log_2\left|1+\frac1x\right|>1.
\]

Сначала:
\[
1=\log_2 2.
\]

Так как $2>1$:
\[
\left|1+\frac1x\right|>2.
\]

А ОДЗ:
\[
\left|1+\frac1x\right|>0
\quad\Longleftrightarrow\quad
1+\frac1x\ne0,
\qquad x\ne0.
\]

Итоговый переход:
\[
\begin{cases}
\left|1+\dfrac1x\right|>2,\\[0.6em]
1+\dfrac1x\ne0,\\[0.4em]
x\ne0.
\end{cases}
\]

\section*{\textcolor{darkaccent}{9. Равносильные переходы для модуля}}
\sectionline

\subsection*{\textcolor{accent}{Модуль больше числа или выражения}}

\[
|F(x)|>G(x)
\quad\Longleftrightarrow\quad
\left[
\begin{array}{l}
F(x)>G(x),\\
F(x)<-G(x).
\end{array}
\right.
\]

Если знак нестрогий:
\[
|F(x)|\ge G(x)
\quad\Longleftrightarrow\quad
\left[
\begin{array}{l}
F(x)\ge G(x),\\
F(x)\le -G(x).
\end{array}
\right.
\]

\subsection*{\textcolor{accent}{Модуль меньше числа или выражения}}

\[
|F(x)|<G(x)
\quad\Longleftrightarrow\quad
\begin{cases}
F(x)<G(x),\\
F(x)>-G(x).
\end{cases}
\]

Если знак нестрогий:
\[
|F(x)|\le G(x)
\quad\Longleftrightarrow\quad
\begin{cases}
F(x)\le G(x),\\
F(x)\ge -G(x).
\end{cases}
\]

\important{Система} означает, что условия выполняются одновременно.  
\important{Совокупность} означает, что достаточно выполнения хотя бы одного условия.

\section*{\textcolor{darkaccent}{10. Когда модуль можно раскрыть только одним способом}}
\sectionline

Иногда ограничение помогает понять знак подмодульного выражения.

\subsection*{\textcolor{accent}{Пример}}

\[
\log_4\frac{4x-5}{|x-1|}\ge\frac12.
\]

Представим:
\[
\frac12=\log_4 2.
\]

Так как $4>1$:
\[
\frac{4x-5}{|x-1|}\ge2.
\]

Это уже гарантирует положительность аргумента логарифма, но само ограничение полезно для анализа:
\[
\frac{4x-5}{|x-1|}>0.
\]

Так как
\[
|x-1|>0 \quad \text{при } x\ne1,
\]
знак дроби определяется числителем:
\[
4x-5>0
\quad\Longleftrightarrow\quad
x>\frac54.
\]

Тогда
\[
x>\frac54>1,
\]
значит
\[
x-1>0,
\qquad
|x-1|=x-1.
\]

Поэтому модуль раскрывается только с плюсом:
\[
\frac{4x-5}{x-1}\ge2.
\]

\important{Вывод:} иногда ограничение кажется лишним, но помогает правильно раскрыть модуль.

\section*{\textcolor{darkaccent}{11. Если модуль стоит не на всём выражении}}
\sectionline

Если модуль входит внутрь выражения, например
\[
x^2-4|x-3|,
\]
то равносильные переходы вида $|F(x)|>G(x)$ напрямую не применяются.

Тогда модуль раскрывают по определению:

\[
|x-3|=
\begin{cases}
x-3, & x-3\ge0,\\
-(x-3), & x-3<0.
\end{cases}
\]

То есть:
\[
|x-3|=
\begin{cases}
x-3, & x\ge3,\\
3-x, & x<3.
\end{cases}
\]

\important{Важно:} знак проверяется не у $x$, а у всего подмодульного выражения.

Например, для $|x-6|$ нужно рассматривать:
\[
x-6\ge0
\quad \text{и} \quad
x-6<0,
\]
а не просто $x\ge0$ и $x<0$.

\section*{\textcolor{darkaccent}{12. Гибридные задачи: логарифмы и степени}}
\sectionline

В сложных неравенствах логарифмы могут появляться внутри показательных выражений.

\subsection*{\textcolor{accent}{Полезная формула}}

\[
a^{\log_a b}=b,
\qquad a>0,\quad a\ne1,\quad b>0.
\]

\subsection*{\textcolor{accent}{Пример другого пути}}

Рассмотрим выражение:
\[
\left(\frac13\right)^{\log_9 F(x)}.
\]

Можно сначала заметить, что
\[
9=\left(\frac13\right)^{-2}.
\]

Но часто проще рассматривать неравенство как показательное.

Например:
\[
\left(\frac13\right)^{\log_9 F(x)}\le1.
\]

Так как
\[
1=\left(\frac13\right)^0,
\]
а основание $\dfrac13$ меньше единицы, получаем:
\[
\log_9 F(x)\ge0.
\]

Дальше:
\[
0=\log_9 1.
\]

Так как $9>1$:
\[
F(x)\ge1.
\]

Ограничение $F(x)>0$ автоматически выполнено, потому что
\[
F(x)\ge1>0.
\]

\section*{\textcolor{darkaccent}{13. Типичные ошибки}}
\sectionline

\begin{enumerate}
    \item Забывать ОДЗ логарифма:
    \[
    f(x)>0.
    \]
    \item Менять знак неравенства при основании $a>1$.
    \item Не менять знак неравенства при основании $0<a<1$.
    \item Писать лишние условия на числовое основание, например $2>0$, $2\ne1$.
    \item Автоматически отбрасывать ограничение, хотя оно не следует из полученного неравенства.
    \item Не использовать полезное ограничение, которое помогает раскрыть модуль.
    \item Для $|x-6|$ ошибочно проверять знак $x$, а не знак $x-6$.
    \item Заменять $|F(x)|>0$ на $F(x)>0$; правильно:
    \[
    |F(x)|>0\Longleftrightarrow F(x)\ne0.
    \]
    \item Путать систему и совокупность при раскрытии модуля.
    \item Забывать знаменатель:
    \[
    \frac{P(x)}{Q(x)} \quad \Rightarrow \quad Q(x)\ne0.
    \]
\end{enumerate}

\section*{\textcolor{darkaccent}{14. Краткая памятка выбора метода}}
\sectionline

\begin{table}[h]
\caption{Что делать в типичных ситуациях}
\centering
\renewcommand{\arraystretch}{1.45}
\begin{tabularx}{\linewidth}{>{\bfseries}p{0.36\linewidth}X}
\toprule
Ситуация & Действие \\
\midrule
Логарифм сравнивается с числом & Представить число как логарифм: $c=\log_a a^c$ \\
Основание $a>1$ & При отбрасывании логарифмов знак сохраняется \\
Основание $0<a<1$ & При отбрасывании логарифмов знак меняется \\
Получили $F(x)\ge2$ & ОДЗ $F(x)>0$ выполнена автоматически \\
Получили $F(x)<1$ & ОДЗ $F(x)>0$ отдельно нужна \\
Есть $|F(x)|$ как весь аргумент логарифма & ОДЗ: $F(x)\ne0$ \\
Есть $|F(x)|>G(x)$ & Раскрыть через совокупность \\
Есть $|F(x)|<G(x)$ & Раскрыть через систему \\
Модуль внутри выражения & Раскрывать по определению через знак подмодульного выражения \\
Логарифм внутри степени & Сначала решить как показательное или использовать $a^{\log_a b}=b$ \\
\bottomrule
\end{tabularx}
\end{table}

\newpage

\section*{\textcolor{darkaccent}{Тренировочный блок}}
\sectionline

Решите неравенства. В каждом задании сначала запишите ОДЗ или объясните, почему часть ограничений выполняется автоматически.

\begin{enumerate}
    \item
    \[
    \log_2(x^2-3x)\ge1.
    \]

    \item
    \[
    \log_{\frac13}\frac{x+2}{x-1}<0.
    \]

    \item
    \[
    \log_4\frac{x^2-4}{x+3}\ge\frac12.
    \]

    \item
    \[
    \log_2(x^2-5x+6)>\log_2(x-2).
    \]

    \item
    \[
    \log_{\frac12}(x^2-4x+3)\le\log_{\frac12}(x-1).
    \]

    \item
    \[
    \log_3\left|\frac{x-2}{x+1}\right|>1.
    \]

    \item
    \[
    \log_4\frac{4x-5}{|x-1|}\ge\frac12.
    \]

    \item
    \[
    \log_{\frac13}\left|\frac{4x+5}{x-1}\right|\le1.
    \]

    \item
    \[
    \log_{\frac12}\left(x^2-4|x-3|\right)\ge -3.
    \]

    \item
    \[
    \left(\frac12\right)^{\log_{\frac14}(x^2-|x-4|)}\ge1.
    \]
\end{enumerate}

\newpage

\section*{\textcolor{darkaccent}{Ответы к тренировочному блоку}}
\sectionline

\begin{longtable}{c>{\raggedright\arraybackslash}p{0.82\linewidth}}
\caption{Ответы}\\
\toprule
№ & Ответ \\
\midrule
\endfirsthead

\toprule
№ & Ответ \\
\midrule
\endhead

1 &
\[
x\in\left(-\infty,\frac{3-\sqrt{17}}2\right]\cup
\left[\frac{3+\sqrt{17}}2,+\infty\right).
\]
\\[0.6em]

2 &
\[
x\in(1,+\infty).
\]
\\[0.6em]

3 &
\[
x\in\left(-3,1-\sqrt{11}\right]\cup
\left[1+\sqrt{11},+\infty\right).
\]
\\[0.6em]

4 &
\[
x\in(4,+\infty).
\]
\\[0.6em]

5 &
\[
x\in[4,+\infty).
\]
\\[0.6em]

6 &
\[
x\in\left(-\frac52,-1\right)\cup
\left(-1,-\frac14\right).
\]
\\[0.6em]

7 &
\[
x\in\left[\frac32,+\infty\right).
\]
\\[0.6em]

8 &
\[
x\in\left(-\infty,-\frac{16}{11}\right]\cup
\left[-\frac{14}{13},1\right)\cup
(1,+\infty).
\]
\\[0.6em]

9 &
\[
x\in\left[-2-2\sqrt6,-6\right)\cup
\left(2,-2+2\sqrt6\right].
\]
\\[0.6em]

10 &
\[
x\in\left(-\infty,\frac{-1-\sqrt{21}}2\right]\cup
\left[\frac{-1+\sqrt{21}}2,+\infty\right).
\]
\\[0.6em]

\bottomrule
\end{longtable}

\end{document}